% \iffalse meta-comment
%
% hit-messages.dtx 是所有面向用户的提示文案
%
% \fi
%
% \section{提示文案}
%
% \changes{v3.2a}{2026/08/09}{面向用户的提示文案从九个模块收进一处}
% 提示文案原先贴着用它的代码写，散在九个 dtx 里。改一句话得先找到它在哪，
% 想知道模板一共会对用户说些什么，只能全仓 grep。
%
% 这里只放 \cs{msg\_new:nnn}，发提示的地方仍旧在各自的模块。分工是：
% 文案在这里改，什么时候发在那边改。
%
% 有四条原先写在宏体里（两条字体、两条数学），正文用的是 \texttt{\#\#1}。
% 搬到顶层之后参数记号还原成 \texttt{\#1}，不然消息里会印出字面的 \texttt{\#1}。
%
% 不在这张表里的只有派发类 \file{hithesis.cls} 的 \texttt{unknown-kind}：
% 那是独立的输出文件，读不到共享层。开发期断言不算提示，就地写，也不进表。
%
% 本文件排在共享层最前，早于任何可能发提示的代码。反复撞上的提示由
% \file{hit-key-engine.dtx} 的 \cs{\_\_hit\_warn\_once:nnn} 按名字去重，只报第一次。
%    \begin{macrocode}
%<*shared-messages>
%    \end{macrocode}
%
% \subsection{旧写法与改名}
%
% 用户写的还是 v3.1e 的写法。模板认，但要告诉他新写法是什么。
% 这一类一律是 warning 不是 error：老文档得先能编出来，用户才有心思改。
%    \begin{macrocode}
\msg_new:nnn { hithesis } { bm-with-unicode-math }
  {
    \iow_char:N \\ usepackage{bm}~跟~unicode-math~不兼容。
    bm~会把~\iow_char:N \\ boldsymbol~换成处理不了扩展~mathchar~的版本，
    公式里一出现希腊字母就报~Extended~mathchar~used~as~mathchar。 \\
    已经把~\iow_char:N \\ bm~与~\iow_char:N \\ boldsymbol~接回~unicode-math~自己的
    \iow_char:N \\ symbf，不用改文档。 \\
    确实要用~bm~本身，请加类选项~math-font=newtx~走~8~位数学字体。
  }
\msg_new:nnn { hithesis } { math-style-with-newtx }
  {
    \iow_char:N \\ documentclass~的~math-style~只对~unicode-math~那几档数学字体
    有效，而现在走的是~fontset-math=newtx~这条~8~位路径。 \\
    newtx~的希腊字母正斜体由~\pkg{newtxmath}~自己决定，math-style~不起作用。 \\
    要国标那套正斜体，请把~fontset-math~换成~termes、xits、stix~等任一档。
  }
\msg_new:nnn { hithesis } { structure-needs-file }
  {
    \iow_char:N \\ hitsetup~里~struct~的~#1~要的是文件名，
    写~auto~没有意义，这一块就不排了。 \\ \\
    auto~是给开关型的部件用的（contents、defense、declarations），
    意思是「让模板决定排不排」。#1~要读一个文件，
    写文件名，或者写~false~明确不排。
  }
\msg_new:nnn { hithesis } { split-option }
  {
    选项~#1~已经拆成两个写法，看你要哪一种。 \\
    \str_case:nn {#1}
      {
        { glue }
          {
            要~v2.0.0~到~v3.1f~那一套版面，连同它手写的伸缩量： \\
            \iow_char:N \\ documentclass[v2-layout={bachelor=true,glue=true}] \\
            只是想让页底排齐，用新版面的写法： \\
            \iow_char:N \\ hitsetup{layout={bottom=flush}} \\
            伸缩量由规范~3.8~节的毫米区间算出来，不用自己填。
          }
        { raggedbottom }
          {
            要~v2.0.0~到~v3.1f~那一套版面： \\
            \iow_char:N \\ documentclass[v2-layout={bachelor=true,bottom=ragged}] \\
            硕博把~bachelor=true~换成~graduate=two。 \\
            只是想调页底，用新版面的写法： \\
            \iow_char:N \\ hitsetup{layout={bottom=ragged}} \\
            页底排齐是~bottom=flush。
          }
      } \\
    旧名这一版照常生效，v4（预计~2027~年发布）中移除。
  }
\msg_new:nnn { hithesis } { renamed-option }
  {
    选项~#1~已改名，请写成~\iow_char:N \\ hitsetup{#2}。 \\
    旧名将在~v4（预计~2027~年发布）中移除。
  }
\msg_new:nnn { hithesis } { renamed-load-option }
  {
    选项~#1~已改名为~#2。 \\
    它决定版面，必须写在~\iow_char:N \\ documentclass~的可选参数里，
    改那一行就行，不要挪进~\iow_char:N \\ hitsetup。 \\
    旧名将在~v4（预计~2027~年发布）中移除。
  }
\msg_new:nnn { hithesis } { renamed-value }
  {
    stage=#1~已废弃，这一份按~#2~处理。 \\
    v3.1e~写成~\iow_char:N \\ documentclass[stage=#1]{hithesisart}；
    v3.2a~起交付物由派发类认，改写成
    \iow_char:N \\ documentclass[kind=#2]{hithesis}。 \\
    直接用~hit-report~类的话，stage~的取值也改叫~#2。 \\
    旧取值将在~v4（预计~2027~年发布）中移除。
  }
\msg_new:nnn { hithesis } { date-not-iso }
  {
    #1~的值不是~ISO~日期： \\
    #2 \\
    原样排出，格式自己负责。 \\
    写成~yyyy、yyyy-mm~或~yyyy-mm-dd。
  }
\msg_new:nnn { hithesis } { date-too-coarse }
  {
    #1~只写到年，这一处要年月： \\
    #2 \\
    有多少排多少。 \\
    补齐写成~yyyy-mm~或~yyyy-mm-dd。
  }
\msg_new:nnn { hithesis } { cover-script-no-font }
  {
    cover-script~写了~font，可这台机器上没装那个字体。 \\
    这一处退回贴图。
  }
\msg_new:nnn { hithesis } { merged-option }
  {
    选项~#1~已并入~#2。 \\
    v4（预计~2027~年发布）中移除。
  }
\msg_new:nnn { hithesis } { renamed-command }
  {
    \iow_char:N \\ #1~已过时，改用~\iow_char:N \\ hitsetup~里的 \\
    #2， \\
    v4（预计~2027~年发布）中移除。
  }
\msg_new:nnn { hithesis } { renamed-environment }
  {
    环境~#1~已改名为~#2。 \\
    旧名将在~v4（预计~2027~年发布）中移除。
  }
\msg_new:nnn { hithesis } { flat-info }
  {
    元信息字段的扁平写法已不推荐，例如~\iow_char:N \\ hitsetup{#1~=~...}， \\
    请改写成~\iow_char:N \\ hitsetup{info~=~{#1~=~...}}。 \\
    扁平写法将在~v4（预计~2027~年发布）中移除。
  }
\msg_new:nnn { hithesis } { style-in-classoptions }
  {
    排版选项~#1~不该写在~\iow_char:N \\ documentclass~的可选参数里， \\
    请改用~\iow_char:N \\ hitsetup。 \\
    在类选项里指定的支持将在~v4（预计~2027~年发布）中移除。
  }
\msg_new:nnn { hithesis } { manual-sides }
  {
    单双面由学位与校区决定，类选项~#1~会被模板覆盖。 \\
    如果确实需要另行指定，请联系模板维护者说明用途。
  }
\msg_new:nnn { hithesis } { bicaption-deprecated }
  {
    \iow_char:N \\ bicaption~已不推荐，英文题注改成~\iow_char:N \\ caption~的
    尾随可选参数： \\
    \iow_char:N \\ caption{中文题}[English~caption]。 \\
    旧写法将在~v4（预计~2027~年发布）中移除。
  }
\msg_new:nnn { hithesis } { deprecated-subfig }
  {
    #2 \\
    请改写成~\iow_char:N \\ subfigure[\iow_char:N \\ bisubcaption{中文}{英文}]%
    {\iow_char:N \\ label{键}图}。 \\
    旧写法将在~v4（预计~2027~年发布）中移除。
  }
\msg_new:nnn { hithesis } { deprecated-ntheorem }
  {
    \iow_char:N \\ #1~是~ntheorem~的命令，模板已改用~thmtools。 \\
    请改写成~\iow_char:N \\ hitsetup{thm~=~{#2~=~...}}。 \\
    兼容写法将在~v4（预计~2027~年发布）中移除。
  }
%    \end{macrocode}
%
% \subsection{已经移除的东西}
%
% 上游宏包换了或者做法废了，旧写法没法再兼容，只能告诉用户改成什么。
%    \begin{macrocode}
\msg_new:nnn { hithesis } { removed-option }
  {
    选项~#1~已废弃，没有对应的新写法，这一句已被忽略。 \\
    #2。 \\
    删掉这一句就行。
  }
\msg_new:nnn { hithesis } { removed-subfig }
  {
    \iow_char:N \\ #1~是~subfigure~宏包的东西，模板已改用~subcaption，这条已经失效。 \\
    请改用~\iow_char:N \\ #2。
  }
\msg_new:nnn { hithesis } { mapped-subfig }
  {
    \iow_char:N \\ #1~是~subfigure~宏包的开关，已自动转成~
    \iow_char:N \\ captionsetup[sub]{#2}。 \\
    建议直接改用新写法，兼容将在~v4（预计~2027~年发布）中移除。
  }
\msg_new:nnn { hithesis } { removed-ntheorem }
  {
    \iow_char:N \\ #1~是~ntheorem~的命令，模板已改用~thmtools，这条已经失效。 \\
    请改用~\iow_char:N \\ #2。
  }
\msg_new:nnn { hithesis } { removed-ntheorem-numbering }
  {
    \iow_char:N \\ #1~是~ntheorem~的命令，模板已改用~thmtools，这条已经失效。 \\
    编号形式请在~\iow_char:N \\ newtheorem~之后自行重定义~\iow_char:N \\ the<环境名>。
  }
%    \end{macrocode}
%
% \subsection{键与取值写错}
%
% 拼错键名，或者取值不在合法集合里。这一类要尽量把合法取值列全，
% 只回显错值等于让用户自己去翻手册。
%    \begin{macrocode}
\msg_new:nnn { hithesis } { unknown-condition }
  {
    \iow_char:N \\ hitsetup~的条件里有不认识的取值~#1，这一组当作不成立。 \\
    条件写成~[kind=thesis]、[campus=shenzhen]~这样，等号右边是选项的取值。
  }
\msg_new:nnn { hithesis } { longbicaption-deprecated }
  {
    \iow_char:N \\ longbionenumcaption~进入废弃期。 \\
    长表格的双语题注改写成~\iow_char:N \\ caption\{中文题\}[英文题]，
    跟浮动体一致；只要中文题就不写后面那个可选参数。 \\
    间距由类里统一给，原来那七个参数里的两个间距不必再自己算。 \\
    旧命令在~v4（预计~2027~年发布）中移除。
  }
\msg_new:nnn { hithesis } { denotation-no-item }
  {
    denotation~里没有~\iow_char:N \\ item，当成旧写法照常排。 \\
    环境名加个星号：\iow_char:N \\ begin\{denotation*\}，
    正文自己写表格，跟以前一样。 \\
    要一行一个符号的话写~\iow_char:N \\ item[符号]~说明，
    第一列宽度用~\iow_char:N \\ begin\{denotation\}[3cm]~改。
  }
\msg_new:nnn { hithesis } { structure-twice }
  {
    structure~里的~#1~在~#2matter~与~#3matter~各写了一次，按后写的~#3matter~排。 \\
    删掉不要的那一处。
  }
\msg_new:nnn { hithesis } { structure-wrong-section }
  {
    structure~里的~#1~不能写在~#2matter~里。 \\
    它只能写在~#3matter。
  }
\msg_new:nnn { hithesis } { unknown-key }
  {
    不认识的键~#1，已忽略。 \\
    检查拼写，键名清单见手册。
  }
\msg_new:nnn { hithesis } { enter-font-unmeasured }
  {
    \token_to_str:N \hitenter ~ 还没有~#1~的自然行高实测值。 \\
    空段行高按「字体自然行高乘倍数」算，自然行高只认实测：
    现有~serif（Times~New~Roman~1.1472）与~songti（宋体~1.29333）。
  }
\msg_new:nnn { hithesis } { bad-option-value }
  { 选项~#1~不接受取值~#2。 }
\msg_new:nnn { hithesis } { measure-clash }
  {
    #1~与~#2~量的是同一段间距，只是度量不同，同一个目标里只能写一个。 \\
    这里按后写的~#2~算。
  }
\msg_new:nnn { hithesis } { key-not-in-language-group }
  {
    键~#1~没有~#2~版本，不能写在~#2~底下。 \\
    它不分语言，请把~#2~去掉；拼错了的话，键名清单见手册。
  }
\msg_new:nnn { hithesis } { font-en-unknown }
  { 选项~font-en~的取值~'#1'~无效。 }
\msg_new:nnn { hithesis } { font-zh-unknown }
  { 选项~font-zh~的取值~'#1'~无效。 }
\msg_new:nnn { hithesis } { math-style-unknown }
  { 选项~math-style~的取值~'#1'~无效，只能是~TeX、ISO~或~GB。 }
\msg_new:nnn { hithesis } { math-font-unknown }
  { 选项~math-font~的取值~'#1'~无效。 }
\msg_new:nnn { hithesis } { msword-no-pitch }
  {
    \iow_char:N \\ hit@msword@layout:n~算不出行距。 \\
    line-pitch~与~lines~至少要给一个。
  }
%    \end{macrocode}
%
% \subsection{排版期发现的}
%
% 编到某一处才暴露的问题，跟用户写的内容有关，不是配置错。
%    \begin{macrocode}
\msg_new:nnn { hithesis } { title-unparsed }
  {
    选项~title-#1~的写法没看懂，整条当成一行处理。 \\
    分两行写成~title={第一行}{第二行}，副标题写成~[副标题]。 \\
    标题本身以花括号开头时（比如~\iow_char:N \\ textbf），外面请再包一层花括号。
  }
\msg_new:nnn { hithesis } { title-too-wide }
  {
    标题第~#1~行排出来宽~#2，放不进题目栏的~#3。 \\
    栏宽是版式定的，改不了；请把~title~的花括号分组挪一挪，让两行都放得下。
  }
\msg_new:nnn { hithesis } { eqdenote-no-item }
  {
    eqdenote~环境里每条注释都要用~\iow_char:N \\item~起头，
    这一处没有找到，排出来的表格会错位。
  }
\msg_new:nnn { hithesis } { bicaption-five-args }
  {
    \iow_char:N \\ bicaption~的五参数写法已不推荐。 \\
    英文图表名请改用~\iow_char:N \\ hitsetup{const={figure-prefix-en=...}}~设置，
    题注写成~\iow_char:N \\ caption{中文题}[English~caption]。 \\
    旧写法将在~v4（预计~2027~年发布）中移除。
  }
\msg_new:nnn { hithesis } { usepackage-in-body }
  {
    \iow_char:N \\ hitusepackage~只能写在导言区，也就是~
    \iow_char:N \\ begin\{document\}~之前。 \\
    \iow_char:N \\ usepackage~本身就有这条限制，这一句已被忽略。
  }
\msg_new:nnn { hithesis } { no-toc-english }
  {
    开题与中期报告只排一份目录，toc~=~\{~english~=~...~\}~不起作用。 \\
    目录的语言跟着类选项~lang~走。
  }
\msg_new:nnn { hithesis } { no-structure }
  {
    这个类没有 #1，主文件里的这一句不起作用。 \\
    前置与后置部分请照旧逐个写 \string\include。
  }
\msg_new:nnn { hithesis } { page-patch-failed }
  {
    没能改写 \string#1。 \\
    版心底留给末行行深的那一段不会全程生效，换页的位置会比 Word 多排一行。 \\
    多半是别的宏包重定义了它，请报 issue。
  }
%    \end{macrocode}
%
% \subsection{参考文献后端}
%
% \changes{v3.2a}{2026/08/16}{参考文献双后端的六条消息}
% \option{bib}/\option{backend} 与两条文献路线相关的消息。biblatex 的一批配置
% （样式、姓名大小写、最多几位作者）在装载那一刻就定死，之后再设不生效，
% 这一类都要说清「什么时候设才算数」。
%    \begin{macrocode}
\msg_new:nnn { hithesis } { bib-backend-in-body }
  {
    backend~=~biber~设晚了：biblatex~必须在~\iow_char:N \\ begin\{document\}~之前
    装载，正文里再设就来不及了。 \\
    请把这一组~bib~=~\{...\}~挪到主文件开头、
    \iow_char:N \\ begin\{document\}~之前。这次设置未生效，仍走传统~BibTeX。
  }
\msg_new:nnn { hithesis } { bib-no-resource }
  {
    biblatex~一线没收到文献库。 \\
    struct~里的~bibliography~要写在~\iow_char:N \\ begin\{document\}~之前，
    才来得及换算成~\iow_char:N \\ addbibresource；
    或者自己在导言区写~\iow_char:N \\ addbibresource\{库.bib\}。 \\
    已经自己~add~过的请忽略这条。
  }
\msg_new:nnn { hithesis } { bib-backend-conflict }
  {
    bib~=~\{~backend~=~bibtex~\}，可是~biblatex~已经装载了。 \\
    两条路线互斥：要用传统~BibTeX~就删掉导言区里装~biblatex~的那一句；
    要用~biblatex~就把~backend~改成~biber，或者干脆不写这个键。
  }
\msg_new:nnn { hithesis } { bib-key-locked }
  {
    bib~里的~#1~设晚了：biblatex~的这项配置在装载那一刻就定死。 \\
    请把~#1~挪到~backend~=~biber~之前，或者跟它写进同一组~bib~=~\{...\}~里。
  }
\msg_new:nnn { hithesis } { bib-style-ignored }
  {
    biblatex~下~\iow_char:N \\ bibliographystyle~无效，这一句已忽略。 \\
    biblatex~一线的样式锁死国标~gb7714-2015，不另选；细调走
    ~bib~=~\{~biber~=~\{...\}~\}~透传。传统一线换~.bst~写
    ~bib~=~\{~bibtex~=~\{~style~=~...~\}~\}。
  }
\msg_new:nnn { hithesis } { package-conflict }
  {
    #2~与已装载的~#1~不能同用。 \\
    参考文献要么全走~biblatex（backend~=~biber），要么全走~natbib~一族，
    两边的引用命令与文献表实现互相覆盖，混装必坏。
  }
%    \end{macrocode}
%
%    \begin{macrocode}
%</shared-messages>
%    \end{macrocode}
%
% \Finale
